\documentclass[11pt]{article}
\usepackage[T1]{fontenc}
\usepackage{textcomp}
\usepackage[margin=0.75in]{geometry}
\usepackage{amsmath,amssymb,amsthm}
\usepackage{array,booktabs}
\usepackage{hyperref}
\setlength{\parindent}{0pt}
\newtheorem{theorem}{Theorem}
\theoremstyle{definition}
\newtheorem{definition}{Definition}
\title{Flow Proof Artifact: Ring-add-cancel}
\author{Generated by \texttt{flow doc proof}}
\date{Flow Proof Book}
\begin{document}
\maketitle
Equal ring elements can be subtracted from both sides.\\[0.5em]
\textbf{Source.} dummit-foote — *Abstract Algebra*, §7.1\\[0.5em]
\bigskip
\noindent{\large \textbf{Derived fact 1.} \textit{a + c = b + c implies a = b} \hfill Ring \cdot addition \cdot right cancellation holds}\par
\smallskip\noindent\textit{Coordinate: Ring \cdot addition \cdot right cancellation holds}\par
\smallskip\noindent\textit{Source: dummit-foote}\par
\smallskip\noindent\textit{Built on: parentheses do not matter, for addition on Ring, zero is the right identity, for addition on Ring}\par
\medskip
\noindent\textbf{Goal.} a + c = b + c implies a = b\par
\noindent$\displaystyle \forall a \in Ring \forall b \in Ring \forall c \in Ring\quad a = b$\par
\medskip
\noindent
\begin{tabular}{@{} >{\raggedright\arraybackslash}p{0.06\textwidth} >{\raggedright\arraybackslash}p{0.47\textwidth} >{\raggedleft\arraybackslash}p{0.41\textwidth} @{}}
\textbf{\#} & \textbf{Proof} & \textbf{Mathematics} \\
\midrule
\textbf{1.} & We prove that right cancellation holds for addition on Ring. & \\[0.45em]
\textbf{2.} & We split into exhaustive cases — the claim must hold in each one. & \\[0.45em]
\textbf{3.} & Case 1 (see step 2): suppose a + c  equals  b + c. & \\[0.45em]
\textbf{4.} & We invoke the derived fact governing addition on Ring: parentheses do not matter, for addition on Ring (instantiated for a, c, -c). & \\[0.45em]
\textbf{5.} & We invoke the derived fact governing addition on Ring: zero is the right identity, for addition on Ring (instantiated for a). & \\[0.45em]
\textbf{6.} & From step 3, step 4, and step 5, this implies a equals b. Together with the other cases (step 3), the goal is discharged. Hence proven. & $\displaystyle a = b$ \\[0.45em]
\bottomrule
\end{tabular}
\label{thm:Ring-addition-right_cancellation_holds}
\par\medskip\hrule
\end{document}
